% Part:sets-functions-relations
% Chapter: sets
% Section: reduction

\documentclass[../../../include/open-logic-section]{subfiles}

\begin{document}

\olfileid[tr]{sfr}{siz}{red}

\olsection{İndirgeme}

\begin{editorial}
  Bu kesim, \olref[nen]{sec}'teki sonuçlara koşut olarak indirgeme yoluyla
  sayılamazlığı ispatlar. \olref[nen-alt]{sec}'teki sonuçlara koşut, biraz
  daha kısa bir alternatif sürüm \olref[red-alt]{sec}'te verilmiştir.
\end{editorial}

Bir köşegenleştirme argümanıyla $\Pow{\PosInt}$'ın !!{nonenumerable}
olduğunu gösterdik. Bütün sonsuz $0$ ve $1$ dizilerinin kümesi
$\Bin^\omega$'nın !!{nonenumerable} olduğuna ilişkin zaten bir ispatımız
vardı. $\Pow{\PosInt}$'ın !!{nonenumerable} olduğunu şöyle başka bir yolla
da ispatlayabiliriz: \emph{$\Pow{\PosInt}$ !!{enumerable} ise
$\Bin^\omega$'nın da !!{enumerable} olduğunu} gösterin. $\Bin^\omega$'nın
!!{enumerable} olmadığını bildiğimize göre $\Pow{\PosInt}$ da olamaz. Buna
bir problemi diğerine \emph{indirgemek} denir; bu durumda $\Bin^\omega$'yı
numaralandırma problemini $\Pow{\PosInt}$'ı numaralandırma problemine
indirgiyoruz. İkincisinin bir çözümü, yani $\Pow{\PosInt}$'ın bir
numaralandırması, birincisinin de bir çözümünü, yani $\Bin^\omega$'nın bir
numaralandırmasını verirdi.

Bir $B$ kümesini numaralandırma problemini bir $A$ kümesini numaralandırma
problemine nasıl indirgeriz? $A$'nın bir numaralandırmasını $B$'nin bir
numaralandırmasına dönüştürmenin yolunu sağlarız. Bunu yapmanın en kolay yolu
bir !!{surjective} $f\colon A\to B$ fonksiyonu tanımlamaktır. $x_1$, $x_2$,
\dots{} $A$'yı numaralandırıyorsa $f(x_1)$, $f(x_2)$, \dots{} da $B$'yi
numaralandırır. Bizim durumumuzda örten bir $f\colon\Pow{\PosInt}\to
\Bin^\omega$ fonksiyonu arıyoruz.

\begin{prob}
Bir !!{injective} $g\colon B\to A$ fonksiyonu varsa ve $B$
!!{nonenumerable} ise $A$'nın da !!{nonenumerable} olduğunu gösterin. Bunu,
$A$'nın bir numaralandırmasını $B$'nin bir numaralandırmasına dönüştürmek
için $g$'yi nasıl kullanabileceğinizi göstererek yapın.
\end{prob}

\begin{proof}[{\olref[nen]{thm:nonenum-pownat}}'ın indirgemeyle ispatı]
$\Pow{\PosInt}$'ın !!{enumerable} olduğunu ve dolayısıyla $Z_1$, $Z_2$,
$Z_3$, \dots{} biçiminde bir numaralandırmasının bulunduğunu varsayalım.

$f\colon\Pow{\PosInt}\to\Bin^\omega$ fonksiyonunu, $f(Z)$ bütün
$n\in\PosInt$ için $s(n)=1$ ancak ve ancak $n\in Z$, aksi hâlde $s(n)=0$
olacak biçimdeki $s$ dizisi olsun diyerek tanımlayın. Bu açıkça bir fonksiyon
tanımlar; çünkü $Z\subseteq\PosInt$ olduğunda herhangi bir $n\in\PosInt$ ya
$Z$'nin bir !!a{element}ıdır ya da değildir. Örneğin pozitif çift sayılar
kümesi $2\PosInt=\{2,4,6,\dots\}$, $010101\dots$ dizisine; boş küme
$0000\dots$ dizisine, $\PosInt$'ın kendisi de $1111\dots$ dizisine eşlenir.

Bu fonksiyon aynı zamanda !!{surjective}dir: Her $0$ ve $1$ dizisi, bir
pozitif tam sayılar kümesine, yani üyeleri dizide $1$ bulunan konumlara
karşılık gelen tam sayılardan oluşan kümeye karşılık gelir. Daha kesin olarak
$s\in\Bin^\omega$ olduğunu varsayın. $Z\subseteq\PosInt$'ı şöyle tanımlayın:
\[
Z = \Setabs{n \in \PosInt}{s(n) = 1}.
\]
O zaman $f$'nin tanımına başvurularak doğrulanabileceği üzere $f(Z)=s$ olur.

Şimdi şu listeyi düşünün:
\[
f(Z_1), f(Z_2), f(Z_3), \dots
\]
$f$ !!{surjective} olduğundan $\Bin^\omega$'nın her üyesi uygun bir argümanda
$f$'nin bir değeri, dolayısıyla listede yer alan bir eleman olmalıdır.
Öyleyse bu liste $\Bin^\omega$'nın tamamını numaralandırmalıdır.

Dolayısıyla $\Pow{\PosInt}$ !!{enumerable} olsaydı $\Bin^\omega$ da
!!{enumerable} olurdu. Ancak $\Bin^\omega$ !!{nonenumerable}dır
(\olref[nen]{thm:nonenum-bin-omega}). O hâlde $\Pow{\PosInt}$
!!{nonenumerable}dır.
\end{proof}

\begin{explain}
İndirgemenin hangi yönde yapıldığı konusunda kafa karışıklığı yaşamak
kolaydır. Örneğin bir !!{surjective} $g\colon\Bin^\omega\to B$ fonksiyonu,
$B$'nin !!{nonenumerable} olduğunu \emph{göstermez}. (Bir $0$ ve $1$ dizisini
ilk !!{element}ına gönderen $g(s)=s(1)$ ile tanımlı $g\colon\Bin^\omega\to
\Bin$ fonksiyonunu düşünün. Bazı diziler $0$, bazıları $1$ ile başladığı için
bu fonksiyon !!{surjective}dir. Fakat $\Bin$ sonludur.) Ayrıca $f$
fonksiyonunun !!{surjective} olması gerektiğine dikkat edin; aksi hâlde
argüman işlemez: $f(x_1)$, $f(x_2)$, \dots{} listesinin $B$'nin bütün
!!{element}larını içereceği güvence altında olmazdı. Örneğin
\[
h(n) = \underbrace{00\dots0}_{\text{$n$ tane $0$}}111\dots
\]
bir $h\colon\PosInt\to\Bin^\omega$ fonksiyonu tanımlar; fakat $\PosInt$
!!{enumerable}dir.
\end{explain}

\begin{prob}
Pozitif tam sayı ikililerinin bütün \emph{kümelerinden} oluşan kümenin bir
indirgeme argümanıyla !!{nonenumerable} olduğunu gösterin.
\end{prob}

\begin{prob}\label{sfr:siz:red:prob:nat-nat}
  Bütün $f\colon\Nat\to\Nat$ fonksiyonlarının kümesi $X$'in bir indirgeme
  argümanıyla !!{nonenumerable} olduğunu gösterin. (İpucu: $X$'ten
  $\Bin^\omega$'ya örten bir fonksiyon verin.)
\end{prob}

\begin{prob}
Doğal sayıların sonsuz dizileri kümesi $\Nat^\omega$'nın bir indirgeme
argümanıyla !!{nonenumerable} olduğunu gösterin.
\end{prob}

\begin{prob}
$P$, pozitif tam sayılar kümesinden $\{0\}$ kümesine giden fonksiyonlar
kümesi; $Q$ da pozitif tam sayılar kümesinden $\{0\}$ kümesine giden
\emph{kısmi} fonksiyonlar kümesi olsun. $P$'nin !!{enumerable}, $Q$'nun ise
öyle olmadığını gösterin. (İpucu: $\Bin^\omega$'yı numaralandırma problemini
$Q$'yu numaralandırma problemine indirgeyin.)
\end{prob}

\begin{prob}
$S$, pozitif tam sayılar kümesinden $\{0,1\}$ kümesine giden bütün
!!{surjective} fonksiyonların kümesi olsun; yani $S$, bütün !!{surjective}
$f\colon\PosInt\to\Bin$ fonksiyonlarından oluşsun. $S$'nin
!!{nonenumerable} olduğunu gösterin.
\end{prob}

\begin{prob}
Bütün gerçel sayılar kümesi $\Real$'in !!{nonenumerable} olduğunu gösterin.
\end{prob}

\end{document}
